\subsection{Добавление и удаление бриков}\label{sec:bricks}

\subsubsection{Distributed-том: добавление брика}

К тому \code{volume2}, состоящему из бриков \code{node1} и \code{node2} (по 5 файлов из 10), добавлен третий брик на \code{node3}. После добавления брика существующие файлы автоматически не перемещаются, поэтому запущена перебалансировка (\code{rebalance}). Затем на клиенте создано ещё 30 файлов, и все 40 файлов оказались распределены между тремя узлами (листинг~\ref{lst:add}).

Количество записей на бриках (13 + 19 + 11 = 43) превышает число файлов (40) на три: это служебные файлы~--- ссылки механизма DHT. При переносе файла на другой брик на прежнем месте остаётся пустой файл-указатель, по которому GlusterFS находит файл. Через точку монтирования клиент их не видит.

\lstfile{5_1_add_brick_distributed.txt}{Добавление брика в distributed-том volume2}\label{lst:add}

\subsubsection{Distributed-том: удаление брика}

Удаление брика выполняется в три этапа: \code{start} (перенос данных с брика на остальные), \code{status} (контроль), \code{commit} (фиксация). После удаления брика \code{node3} все 40 файлов остались доступны на клиенте, а на самом брике файлов не осталось (листинг~\ref{lst:remove}). Предупреждение об опции \code{cluster.force-migration} означает, что файлы, в которые ведётся запись во время миграции, могут быть перенесены некорректно, поэтому при удалении брика запись в том желательно остановить.

\lstfile{5_2_remove_brick_distributed.txt}{Удаление брика из distributed-тома volume2}\label{lst:remove}

\subsubsection{Replica-том: изменение числа реплик}

В replica-томе \code{volume1} число копий увеличено с двух до трёх добавлением брика на \code{node3}. Служба самолечения (self-heal) скопировала на новый брик файл \code{test0.txt}. Затем брик удалён, и том возвращён к конфигурации replica~2 (листинг~\ref{lst:replica}).

\lstfile{5_3_replica_add_remove.txt}{Изменение числа реплик тома volume1}\label{lst:replica}
